% =============================================================================
% ch07_oursystem.tex — Chapter 7: Visual-Inertial Odometry and Visual SLAM | NavigatorCrew
% Compiler: XeLaTeX | Language: Hebrew primary, English technical terms via \en{}
% =============================================================================

\selectlanguage{hebrew}

\section{אודומטרייה חזותית-אינרציאלית ו-\en{SLAM} חזותי}
\label{sec:oursystem}

% ── 7.1 ORB-SLAM3 ───────────────────────────────────────────────────────────
\subsection{\en{ORB-SLAM3}: מצב-האמנות ב-\en{SLAM} חזותי}
\label{subsec:orb_slam3}

\en{ORB-SLAM3} מייצג את מצב-האמנות הנוכחי ב-\en{SLAM} חזותי
מבוסס תכונות. זוהי המערכת הראשונה המסוגלת לפעול במצבים מונוקולרי,
סטריאו ו-\en{RGB-D} עם \en{Fusion} חזותי-אינרציאלי, תוך תמיכה
בדגמי מצלמת חור-סיכה (\en{Pin-Hole}) ועין-דג (\en{Fisheye}).
החדשנות המרכזית היא מסגרת ה-\en{Atlas}, המנהלת מבנה נתונים
רב-מפתי המאפשר התאוששות חזקה מכשלי מעקב על ידי יצירת מפות
חדשות ומיזוגן מאוחר יותר כאשר מתגלות סגירות לולאה
\cite{MurArtal2015ORB}.

על מערך הנתונים \en{EuRoC MAV}, \en{ORB-SLAM3} משיג \en{ATE}
של 0.03 מטר לרצף \en{V1_01} במצב סטריאו-אינרציאלי, תוצאה טובה
פי 2.6 מ-\en{VINS-Mono}. האלגוריתם מבצע \en{Bundle Adjustment}
חזותי-אינרציאלי המשלב שגיאות תחזית מחדש (\en{Reprojection Error})
עם שגיאות \en{IMU Preintegration}, תוך שימוש בפונקציות עלות
\en{Robust} (\en{Huber Loss}) לטיפול במדידות חריגות.

\begin{equation}
    \mathbf{X}^* = \arg\min_{\mathbf{X}} \sum_{i} \rho_{\delta}\left( \| \mathbf{e}_{\text{vis},i} \|_{\Omega_{\text{vis}}} \right)
    + \sum_{j} \rho_{\delta}\left( \| \mathbf{e}_{\text{imu},j} \|_{\Omega_{\text{imu}}} \right)
    \label{eq:orb_slam3_ba}
\end{equation}

משוואה~\ref{eq:orb_slam3_ba} מציגה את פונקציית העלות של
\en{Bundle Adjustment} חזותי-אינרציאלי, כאשר $\mathbf{e}_{\text{vis}}$
היא שגיאת התחזית מחדש ו-$\mathbf{e}_{\text{imu}}$ היא שגיאת
\en{IMU Preintegration}.

% ── 7.2 DSO and SVO ─────────────────────────────────────────────────────────
\subsection{\en{DSO} ו-\en{SVO}: גישות ישירות לעומת מבוססות תכונות}
\label{subsec:dso_svo}

\en{Direct Sparse Odometry} (\en{DSO}) נוקט גישה שונה מהותית:
במקום לחלץ ולהתאים תכונות, \en{DSO} ממזער את השגיאה הפוטומטרית
על קבוצה דלילה של פיקסלים, תוך אופטימיזציה משותפת של תנוחת
המצלמה, העומק ההפוך ופרמטרי בהירות אפינית. \en{DSO} משיג דיוק
מעולה בסביבות עשירות בטקסטורה, אך נכשל בסביבות דלות טקסטורה
או בתנאי טשטוש תנועה.

\en{Semi-Direct Visual Odometry} (\en{SVO}) משלב שיטות ישירות
ומבוססות תכונות: הוא משתמש ביישור ישיר (\en{Direct Alignment})
לאומדן תנועה בין-פריימית ושיטות מבוססות תכונות למיפוי וסגירת
לולאה. \en{SVO 2.0} משיג \en{ATE} של 0.09 מטר על \en{EuRoC V1_01}
עם עומס מעבד נמוך של 28%, מה שהופך אותו למתאים במיוחד
לפלטפורמות מוגבלות משאבים כמו רחפנים.

\begin{equation}
    \mathbf{T}^* = \arg\min_{\mathbf{T}} \sum_{i \in \mathcal{P}} \left\| I_1(\mathbf{p}_i) - I_2(\pi(\mathbf{T} \cdot \pi^{-1}(\mathbf{p}_i, d_i))) \right\|_{\sigma}^2
    \label{eq:dso_photometric}
\end{equation}

משוואה~\ref{eq:dso_photometric} מגדירה את השגיאה הפוטומטרית
ב-\en{DSO}, כאשר $I_1, I_2$ הן עוצמות הפיקסלים בשתי תמונות
עוקבות, $\pi$ היא פונקציית ההיטל של המצלמה, ו-$d_i$ הוא העומק
ההפוך של פיקסל $i$.

\begin{equation}
    \mathbf{T}^* = \arg\min_{\mathbf{T}} \sum_{i} \left\| I_{\text{ref}}(\mathbf{p}_i) - I_{\text{cur}}(\mathbf{w}(\mathbf{T}, \mathbf{p}_i, d_i)) \right\|^2
    \label{eq:svo_direct}
\end{equation}

משוואה~\ref{eq:svo_direct} מציגה את יישור התמונה הישיר
ב-\en{SVO}, כאשר $\mathbf{w}$ היא פונקציית העיוות
(\en{Warp function}) הממפה פיקסלים מתמונת הייחוס לתמונה הנוכחית.

% ── 7.3 Comparative Analysis ────────────────────────────────────────────────
\subsection{ניתוח השוואתי ומגבלות}
\label{subsec:visual_comparison}

השוואת ביצועים על \en{EuRoC} ו-\en{TUM VI} מראה כי
\en{ORB-SLAM3} מוביל בדיוק בכל הרצפים, עם \en{ATE} של
0.03-0.08 מטר לעומת 0.08-0.24 מטר למתחרים. עם זאת,
\en{ORB-SLAM3} דורש עומס מעבד גבוה יותר (45%) לעומת
\en{SVO} (28%). \en{DSO} סובל מביצועים ירודים בסביבות
דלות טקסטורה, בעוד ש-\en{VINS-Mono} מתקשה בתנועה סיבובית
מהירה.

כשלים מרכזיים למערכות \en{SLAM} חזותי כוללים: תנועה סיבובית
טהורה הגורמת לאובדן מעקב תכונות, סביבות דלות טקסטורה (קירות
לבנים, מנהרות), שינויי תאורה מהירים, ועצמים דינמיים המפרים
את הנחת העולם הסטטי. \en{Fusion} חזותי-אינרציאלי מפחית חלק
מבעיות אלה על ידי אספקת קנה מידה מטרי ואומדני זווית מיושרי
כבידה מה-\en{IMU}, אך שגיאות \en{IMU Preintegration} מצטברות
בתקופות עמידה ממושכות.

\begin{table}[H]
\centering
\begin{tabular}{lccc}
\toprule
אלגוריתם & \en{ATE} [מ׳] & עומס מעבד [%] & חסינות לחושך \\
\midrule
\en{ORB-SLAM3} & 0.03-0.08 & 45 & נמוכה \\
\en{SVO 2.0} & 0.09-0.15 & 28 & נמוכה \\
\en{DSO} & 0.12-0.24 & 35 & נמוכה מאוד \\
\en{VINS-Mono} & 0.08-0.18 & 40 & נמוכה \\
\bottomrule
\end{tabular}
\caption{השוואת ביצועי מערכות \en{SLAM} חזותי על \en{EuRoC MAV}.}
\label{tab:visual_slam_benchmark}
\end{table}

\begin{figure}[H]
    \centering
    \includegraphics[width=0.92\columnwidth]{figures/fig_visual_slam_comparison.png}
    \caption{השוואת מסלולי \en{SLAM} חזותי על רצף \en{EuRoC V1_01}.
             \en{ORB-SLAM3} (כחול) קרוב ביותר ל-\en{Ground Truth} (שחור).}
    \label{fig:visual_slam_comparison}
\end{figure}

% ── 7.4 Summary ─────────────────────────────────────────────────────────────
\subsection{סיכום}
\label{subsec:visual_summary}

בפרק זה סקרנו את הגישות המרכזיות ל-\en{SLAM} חזותי ואודומטרייה
חזותית-אינרציאלית. \en{ORB-SLAM3} מציג את הביצועים הטובים ביותר
במונחי דיוק, אך בעלות חישובית גבוהה. \en{SVO} מהווה חלופה
יעילה לפלטפורמות מוגבלות משאבים. בפרק הבא נתאר את מערכת
ה-\en{Fusion} הרב-חיישני המשלבת את כל המודאליות שתוארו.